\glsresetall

\chapter{Katamaran}
% \begin{itemize}
%     \item what is katamaran
%     \begin{itemize}
%         \item using Universal contracts
%         \item applied to \gls{pmp}
%     \end{itemize}
%     \item we want to extend to include \gls{smepmp}
% \end{itemize}

% Katamaran is a tool to formally verify \textmu{}sail programs such as ISA definitions, it does this by verifing the code against a \gls{uc}.
% A \gls{uc} are specifications of the intended behaviors of the \textmu{}sail program to verify.
% In the case of ISA definitions these \glspl{uc} can be the security guarantees for the ISA.
% Katamaran can then mathematically prove whether the \textmu{}sail code adheres to the \gls{uc} instead of needing to have a whole bunch of tests to attempt to cover the whole codebase.~\cite{universal_contract}

Katamaran is a framework for formally verifying \textmu{}Sail programs, including ISA definitions, through the use of \glspl{uc}.
\Glspl{uc} formally specify the intended behaviour of a \textmu{}Sail program using preconditions and postconditions.
For ISA models, these contracts can be used to express architectural security guarantees.
Katamaran then proves that the \textmu{}Sail implementation satisfies these contracts, enabling mathematical guarantees about correctness and security beyond what can be achieved through testing alone.~\cite{universal_contract}

The Katamaran paper includes a case study for the RISC-V \gls{pmp} extension.
This case study formalizes the PMP specification in \textmu{}Sail and defines a \gls{uc} capturing the security properties related to \gls{pmp}'s memory isolation and access control.

% TODO
This \textmu{}Sail ISA model is derived from the official RISC-V Sail implementation for \gls{pmp}~\cite{risk-v_sail}.
Sail is a language designed to describe ISA semantics and is used by RISC-V to define it's standard.
These models are translated into \textmu{}Sail to be used in Katamaran as a representative model closely aligned with the official RISC-V specification.

% The model is however limited to just a proof of concept instead of a full \gls{pmp} model, it only allows for only 2 \gls{pmp} configurations instead of 16 or 64 and only allows the \gls{tor} addressing mode.% TODO see GPT output
The current PMP model is intentionally limited and primarily serves as a proof of concept rather than a complete PMP implementation.
It supports only two \gls{pmp} configurations instead of the 16 or 64 entries supported by the RISC-V specification, and only implements the \gls{tor} addressing mode.

The model also does not include the \gls{smepmp} extension.
This thesis extends the existing Katamaran \gls{pmp} case study with support for \gls{smepmp} by extending the ISA model and the associated formal verification.%adding the required architectural state, updating the \gls{pmp} semantics and extending the associated proofs.
